% 供「三、模型假设」增补问题三条。已有 Q1/Q2 公共假设不再重复。

\item 问题三的普通购电账单锚定每日 0:00 原计划 \(g^0\)，最终有效承诺为时段到来前的 \(g^F\)；6:00、12:00、18:00 的中间版本不对未执行电量重复收费。这是避免重复计价的本文解释，不是题面逐字唯一的市场合同。

\item 主模型在 0:00 及随后更新点规划未来时段时，负荷只用日期严格早于当天的最近至多 4 个同星期日均值；真实负荷仅在当前 10 分钟执行步进入能量平衡。当天完整实际负荷曲线只作为理想化信息对照，不作为可执行主方案。

\item 附件 3 的小时光伏预报映射到 10 分钟时，以更新时点已观测实际光伏为锚点并作线性插值。该映射是数据粒度处理，不称为附件规定的唯一细分规则。

\item 6:00、12:00、18:00 只能改写其后尚未执行的普通购电承诺；已发生的购电、SOC 与充放电不得被新预测改写。储能与紧急购电只基于已观测信息补救，不能替代本可在最近更新点完成的正常购电调整。

\item 日末 SOC 不强制回到日初。主核算拟取年末 1200 kWh 作为统一边界，以便与问题二已签收路径比较库存口径；6000 kWh 仅作库存中性敏感性。若问题二政策一致全年结果的年末库存不再接近下界，则改为“本文核算边界”，不把 1200 kWh 写成题面规定或社区运行规律。

\item 下一日到达成本 \(V_{d+1}\) 只引导今日库存决策，不计入当日购电账单，也不提前锁定次日计划。
